\documentclass[11pt]{article}
\usepackage[utf8]{inputenc}
\usepackage[T1]{fontenc}
\usepackage{geometry}
\usepackage{graphicx}
\usepackage{titlesec}
\usepackage{enumitem}
\usepackage{hyperref}
\usepackage{listings}
\usepackage{xcolor}

\geometry{a4paper, margin=2.5cm}

\definecolor{codegray}{rgb}{0.9,0.9,0.9}
\lstset{backgroundcolor=\color{codegray}, frame=single, basicstyle=\ttfamily\small, breaklines=true}

\title{\textbf{Credit Platform}\\Rapport Technique Détaillé - Version 2.0}
\author{Équipe de Développement}
\date{\today}

\begin{document}

\maketitle
\tableofcontents
\newpage

\section{1. Rappel du Concept & Problématique}
\subsection{Le Défi}
Les institutions financières reçoivent des milliers de demandes de crédit non structurées via des canaux hétérogènes (Email, WhatsApp, Portails). Le traitement manuel entraîne :
\begin{itemize}
    \item Temps de réponse longs (>48h)
    \item Erreurs de saisie
    \item Manque de traçabilité et de contexte client
\end{itemize}

\subsection{La Solution : Credit Platform Automation}
Une plateforme centrale d'ingestion et d'analyse intelligente qui transforme des messages bruts en dossiers structurés prêts pour décision, réduisant le temps de traitement à quelques minutes grâce à l'IA vectorielle.

\section{2. Architecture Technique Détaillée}

\subsection{Stack Technologique}
\begin{itemize}
    \item \textbf{Backend} : Python FastAPI (Performance, Async)
    \item \textbf{Base de Données} : MongoDB (Flexibilité schéma JSON)
    \item \textbf{Moteur IA} : Qdrant Cloud (Vector Search Engine)
    \item \textbf{NLP} : spaCy (Extraction d'entités), sentence-transformers (Embeddings)
\end{itemize}

\subsection{Intégration Qdrant (Cœur IA)}
Contrairement aux classifieurs classiques, nous utilisons la \textbf{Recherche par Similarité Vectorielle} (RAG simplifié).
\begin{itemize}
    \item \textbf{Modèle d'embedding} : \texttt{all-MiniLM-L6-v2} (rapide et précis pour le sémantique)
    \item \textbf{Collection} : \texttt{email\_memory}
    \item \textbf{Configuration} :
    \begin{itemize}
        \item Dimension : 384
        \item Distance : Cosine (mesure l'angle sémantique entre deux textes)
    \end{itemize}
\end{itemize}

\section{3. Pipeline de Données (Workflow)}

\subsection{Phase 1 : Ingestion Multi-Sources}
Récupération asynchrone des messages.
\begin{lstlisting}
{
  "source": "gmail",
  "sender": "client@email.com",
  "content": "Je veux 50k pour une voiture...",
  "timestamp": "2026-01-26T10:00:00"
}
\end{lstlisting}

\subsection{Phase 2 : Traitement NLP & Enrichissement}
Le script \texttt{process\_messages.py} nettoie le texte et extrait les entités clés :
\begin{itemize}
    \item \textbf{Montant} : Détection contextuelle ("50k", "50 000 dt")
    \item \textbf{Type} : Déduit par mots-clés initiaux (ex: "voiture" -> AUTO)
    \item \textbf{Client} : Extraction du CIN, Tél, Email via Regex/NER
\end{itemize}

\subsection{Phase 3 : Intelligence Vectorielle (Future)}
C'est ici que l'IA intervient pour comprendre l'intention réelle.
\begin{enumerate}
    \item Le message nettoyé est converti en vecteur (384 float).
    \item Qdrant cherche les 5 vecteurs les plus proches dans le dataset de référence.
    \item Le système déduit l'intention (Ex: "Demande Crédit") et le ton (Ex: "Urgent") par vote majoritaire des voisins.
\end{enumerate}

\section{4. État d'Avancement (Timeline)}

\subsection{Travaux Réalisés (Work Done)}
\begin{enumerate}
    \item \textbf{Infrastructure} :
    \begin{itemize}
        \item Serveur FastAPI opérationnel
        \item Connexion MongoDB Atlas validée
        \item Gestion des variables d'environnement (.env) sécurisée
    \end{itemize}
    \item \textbf{Ingestion} :
    \begin{itemize}
        \item Module Gmail (OAuth2) : Fonctionnel
        \item Module WhatsApp (Webhook Meta) : Prêt
        \item Module Banque (API simulation) : Prêt
    \end{itemize}
    \item \textbf{Traitement de Données} :
    \begin{itemize}
        \item Pipeline de nettoyage NLP implémenté
        \item Stockage structuré des messages bruts et traités
    \end{itemize}
\end{enumerate}

\subsection{Travaux à Venir (To Do)}
L'accent est mis sur l'activation du moteur IA :
\begin{enumerate}
    \item \textbf{Setup Qdrant (En cours)} :
    \begin{itemize}
        \item Configuration du cluster Cloud (Effectué)
        \item Création des collections (Effectué)
    \end{itemize}
    \item \textbf{Dataset Synthétique} :
    \begin{itemize}
        \item Génération de 100+ emails de référence annotés (Intent + Tone)
        \item Validation manuelle de la qualité
    \end{itemize}
    \item \textbf{Vectorisation} :
    \begin{itemize}
        \item Implémentation du script d'embedding (nécessite résolution dépendances PyTorch)
        \item Indexation massive dans Qdrant
    \end{itemize}
    \item \textbf{Moteur de Décision} :
    \begin{itemize}
        \item Logique métier : Si "Urgent" et "Risque élevé" -> Alerte Analyste Senior
    \end{itemize}
\end{enumerate}

\end{document}
